%% sections/05-conclusion.tex
%% Section 5: Conclusion and Outlook (~2 pages)

\section{Conclusion and Outlook}\label{sec:conclusion}

\todo{Content to fill in:
(a) Summary table: the strong vs.\ weak spike contrast, three rows:
    eigenvalue limit, eigenvector overlap, fluctuation scale.

    | Spike strength | Eigenvalue limit  | Overlap with \(\vect e_1\) | Fluctuation |
    |----------------|-------------------|----------------------------|-------------|
    | Strong \(\alpha > 1 + \sqrt y\) | \(\Psi(\alpha)\) | \(\rho(\alpha, y) > 0\) | Gaussian, \(\sqrt n\) |
    | Weak \(1 < \alpha \le 1 + \sqrt y\) | \((1 + \sqrt y)^2\) | 0 | Edge / Tracy-Widom, \(p^{-2/3}\) |

(b) One-paragraph synthesis: the spike phase transition is the cleanest
    quantitative example of when high-dimensional inference recovers
    structure and when it does not. It explains, in a controlled toy
    model, why methods like PCA and spectral clustering succeed only
    when signal exceeds a threshold determined by the aspect ratio.

(c) Outlook: open questions to flag.
    - The critical case \(\alpha = 1 + \sqrt y\). The current literature
      does not yet contain a clean limit theorem; cite open work.
    - Exact variance \(\sigma_{\alpha, y}^2\) in the strong-spike CLT.
    - Beyond Johnstone's case: general spiked models with non-flat bulk
      (cite \citet{BaikSilverstein2006, BenaychGeorgesNadakuditi2011}).
    - Rigorous spectral clustering derivations beyond the toy model;
      cite \citet{CouilletBenaych2016}.

(d) Closing sentence: the goal of this thesis was to develop the
    foundational picture from one end of the dimension spectrum to the
    other; the open questions above point to where the field is going.}

\clearpage
